%! Linear Combinations of Vectors — Unit 1, Lesson 1.1 — Slides (Beamer)
\documentclass[aspectratio=169, 11pt]{beamer}
\usepackage{linalg-beamer}   % provides \burgundyheader{} and \sectionlabel[color]{}
\newcommand{\vv}{\mathbf{v}}
\newcommand{\ww}{\mathbf{w}}

\begin{document}

% TITLE SLIDE (CourseName is not defined in beamer, so the name is literal)
{
\setbeamercolor{background canvas}{bg=burgundy}
\begin{frame}[plain]
  \vspace{2.8cm}\hspace{0.55cm}%
  \begin{minipage}{0.9\textwidth}
    {\color{goldacc}\bfseries\small LESSON 1.1}\\[0.5cm]
    {\color{white}\bfseries\LARGE Linear Combinations of Vectors}\\[0.3cm]
    {\color{white}\bfseries\large Scaling, adding, and what you can reach}\\[1.0cm]
    {\color{blushmid}\small Linear Algebra~$\cdot$~Unit 1: Vectors and Matrices}
  \end{minipage}
\end{frame}
}

% HOOK
\begin{frame}[t, plain]
  \burgundyheader{Two moves, how far can you go?}
  \sectionlabel{Hook}
  \begin{columns}[T]
    \begin{column}{0.58\textwidth}
      A robot at the origin has only two moves it can scale and repeat:
      $\vv=\begin{bsmallmatrix} 2\\1 \end{bsmallmatrix}$ and
      $\ww=\begin{bsmallmatrix} 1\\3 \end{bsmallmatrix}$.\\[8pt]
      \textbf{Question.} Using $c$ copies of $\vv$ and $d$ copies of $\ww$, can it land on
      $(5,5)$? On \emph{any} target?
    \end{column}
    \begin{column}{0.40\textwidth}
      \centering
      \begin{tikzpicture}[scale=0.5]
        \draw[step=1, black!15, thin] (-0.3,-0.3) grid (5.3,5.3);
        \draw[->, thick, black!60] (0,0)--(5.4,0);
        \draw[->, thick, black!60] (0,0)--(0,5.4);
        \draw[->, very thick, burgundy] (0,0)--(2,1) node[below right]{$\vv$};
        \draw[->, very thick, royalblue] (0,0)--(1,3) node[above left]{$\ww$};
        \fill[black!70] (5,5) circle (3pt) node[above right]{\scriptsize $(5,5)$};
      \end{tikzpicture}
    \end{column}
  \end{columns}
\end{frame}

% WHAT IS A LINEAR COMBINATION
\begin{frame}[t, plain]
  \burgundyheader{What is a linear combination?}
  \sectionlabel{Definition}
  A \textbf{linear combination} scales each vector, then adds:
  \[
    c\,\vv + d\,\ww .
  \]
  The scalars $c,d$ are the \textbf{coefficients} (weights).
  \begin{itemize}
    \item Example: $2\vv+\ww=\begin{bsmallmatrix} 4\\2 \end{bsmallmatrix}+\begin{bsmallmatrix} 1\\3 \end{bsmallmatrix}=\begin{bsmallmatrix} 5\\5 \end{bsmallmatrix}$.
    \item Familiar cases: the \textbf{sum} $\vv+\ww$, the \textbf{difference} $\vv-\ww$, a lone \textbf{multiple} $c\,\vv$ ($d=0$).
  \end{itemize}
\end{frame}

% ONE VECTOR = A LINE
\begin{frame}[t, plain]
  \burgundyheader{One vector fills a line}
  \sectionlabel{Span}
  \begin{columns}[T]
    \begin{column}{0.55\textwidth}
      Scale a single $\vv$ by \emph{every} number $c$: all the points $c\,\vv$ lie on one
      \textbf{line through the origin}.\\[8pt]
      That set of reachable points is the \textbf{span} of $\vv$.
    \end{column}
    \begin{column}{0.42\textwidth}
      \centering
      \begin{tikzpicture}[scale=0.5]
        \draw[step=1, black!15, thin] (-2.3,-1.3) grid (4.3,2.3);
        \draw[->, thick, black!60] (-2.4,0)--(4.4,0);
        \draw[->, thick, black!60] (0,-1.4)--(0,2.4);
        \draw[burgundy!55, very thick] (-2.2,-1.1)--(4.2,2.1);
        \draw[->, very thick, burgundy] (0,0)--(2,1) node[above left]{$\vv$};
        \fill[burgundy] (4,2) circle (2.5pt);
        \fill[burgundy] (-2,-1) circle (2.5pt);
      \end{tikzpicture}
    \end{column}
  \end{columns}
\end{frame}

% TWO VECTORS = THE PLANE
\begin{frame}[t, plain]
  \burgundyheader{Two vectors fill the plane (usually)}
  \sectionlabel[goldacc]{Span}
  Let both weights vary in $c\,\vv+d\,\ww$:
  \begin{itemize}
    \item If $\vv$ and $\ww$ are \textbf{not parallel}, combinations reach \textbf{every} point in the plane.
    \item The basis $\mathbf{i}=\begin{bsmallmatrix} 1\\0 \end{bsmallmatrix},\ \mathbf{j}=\begin{bsmallmatrix} 0\\1 \end{bsmallmatrix}$:
          every $\begin{bsmallmatrix} x\\y \end{bsmallmatrix}=x\,\mathbf{i}+y\,\mathbf{j}$ --- coordinates \emph{are} the weights.
  \end{itemize}
  \begin{block}{The exception}
    If $\vv$ and $\ww$ are \textbf{parallel} (one a multiple of the other), their combinations
    span only a \textbf{line}.
  \end{block}
\end{frame}

% WRAP / PREVIEW
\begin{frame}[t, plain]
  \burgundyheader{Wrap-up and what's next}
  \sectionlabel{Connections}
  \textbf{Today:} linear combinations $c\,\vv+d\,\ww$; one vector spans a line, two
  non-parallel vectors span the plane; finding the weights for a target.\\[8pt]
  \textbf{Next --- Lesson 1.2:} the \textbf{dot product} --- a single number measuring how much
  two vectors point the same way, and how it recovers length.
\end{frame}

\end{document}
